\documentclass[11pt]{article}
\usepackage{amsmath,amssymb,amsthm}
\usepackage{geometry}
\usepackage{booktabs}
\usepackage{graphicx}
\usepackage{hyperref}
\usepackage{enumerate}
\usepackage{listings}
\usepackage{xcolor}
\usepackage{array}
\usepackage{enumitem}

\geometry{letterpaper, margin=1in}

\lstset{
  basicstyle=\ttfamily\small,
  keywordstyle=\color{blue},
  commentstyle=\color{gray},
  stringstyle=\color{red},
  breaklines=true,
  columns=fullflexible,
  keepspaces=true,
  language=C++
}

\newcommand{\file}[1]{\texttt{#1}}
\newcommand{\type}[1]{\texttt{#1}}
\newcommand{\field}[1]{\texttt{#1}}
\newcommand{\op}[1]{\texttt{#1}}
\newcommand{\ns}[1]{\texttt{#1}}

\title{\textbf{Formation Engine Methodology} \\
\Large Coarse-Grained and Multi-Scale Layer \\[6pt]
\normalsize Technical Architecture Document}
\author{Formation Engine Development Team}
\date{Version 0.9 --- Phases 1--8 Complete}

\begin{document}
\maketitle
\tableofcontents
\newpage

% =============================================================================
\section{Executive Summary}
% =============================================================================

This document specifies the architecture and integration plan for the
\textbf{coarse-grained (CG) and multi-scale layer} of the formation engine.

The atomistic layer (documented in Sections~1--9 of the existing methodology)
provides a validated, deterministic simulation kernel operating on systems of
2--10,000 atoms with explicit positions, velocities, and forces.  It produces
coordinates, energies, forces, and basic validation on discrete systems.

The coarse-grained layer extends this foundation to larger length and time
scales by introducing:

\begin{enumerate}
  \item A mapping framework that converts atomistic outputs into reduced
        bead representations
  \item Statistical mechanics wrappers that produce ensemble-averaged
        observables rather than single configurations
  \item Emergent-behavior analysis that extracts collective structural
        properties from sampled trajectories
  \item Multi-scale bridges that propagate information consistently between
        atomistic, coarse-grained, and continuum-like descriptions
  \item Validation and benchmarking infrastructure that co-evolves with
        every new capability
\end{enumerate}

The atomistic layer is no longer merely a simulator.  It becomes a
\textbf{teacher model} for the coarse-grained layer.

\subsection{Terminology Note}

Per project terminology rules, this document uses ``atomistic'' for all
particle-resolved descriptions and ``coarse-grained'' (CG) for reduced
representations.  The term ``multi-scale'' refers to the coupling between
resolution levels.  The scale regime beyond $10{,}000$ atoms that was
previously labeled in \S10 is referred to here as the
\textbf{coarse-grained regime}.

% =============================================================================
\section{Atomistic Foundation: Current Capabilities}
% =============================================================================

The existing atomistic kernel provides the reference data that all
coarse-grained methods consume.  This section summarizes what is operational
and where the scaling limits occur.

\subsection{Operational Capabilities}

The formation engine currently supports:

\begin{center}
\begin{tabular}{lll}
\toprule
\textbf{Capability} & \textbf{Implementation} & \textbf{Status} \\
\midrule
State representation & \file{atomistic/core/state.hpp} & Complete \\
LJ 12-6 + Coulomb forces & \file{atomistic/models/lj\_coulomb.cpp} & Complete \\
Bonded interactions & \file{atomistic/models/bonded.cpp} & Complete \\
Composite model & \file{atomistic/models/composite.cpp} & Complete \\
FIRE minimization & \file{atomistic/integrators/fire.cpp} & Complete \\
Velocity Verlet (NVE) & \file{atomistic/integrators/velocity\_verlet.cpp} & Complete \\
Langevin dynamics (NVT) & \file{atomistic/integrators/velocity\_verlet.cpp} & Complete \\
Maxwell--Boltzmann init & \file{atomistic/core/maxwell\_boltzmann.cpp} & Complete \\
PBC (orthogonal box) & \file{atomistic/core/state.hpp} & Complete \\
Supercell construction & \file{atomistic/crystal/supercell.cpp} & Complete \\
Crystal metrics & \file{atomistic/crystal/crystal\_metrics.cpp} & Complete \\
PMF extraction & \file{atomistic/analysis/pmf.cpp} & Complete \\
Structure classification & \file{atomistic/classify/cluster.cpp} & Complete \\
Welford online statistics & \file{atomistic/core/statistics.hpp} & Complete \\
Stationarity gating & \file{atomistic/core/statistics.hpp} & Complete \\
Kabsch alignment / RMSD & \file{atomistic/core/alignment.cpp} & Complete \\
XYZ parsing / compilation & \file{atomistic/parsers/}, \file{compilers/} & Complete \\
\bottomrule
\end{tabular}
\end{center}

\subsection{Validated Outputs}

The atomistic kernel produces the following validated outputs that serve as
inputs to the coarse-grained layer:

\begin{itemize}
  \item \textbf{Equilibrated structures}: Local energy minima found via
        MD exploration + FIRE quenching (\S6)
  \item \textbf{Trajectories}: Time-ordered sequences of states from NVE or
        NVT dynamics (\S5)
  \item \textbf{Energy decomposition}: Per-term energy ledger
        $\mathcal{E} = (U_\text{bond}, U_\text{angle}, U_\text{tors},
        U_\text{vdW}, U_\text{Coul}, U_\text{ext}, U_\text{pol})$ (\S3)
  \item \textbf{Forces}: Per-atom force vectors $\mathbf{F}_i$ (\S3)
  \item \textbf{Radial distribution functions}: $g(r)$ from trajectory
        analysis (\file{analysis/pmf.cpp})
  \item \textbf{Potentials of mean force}: $W(r) = -k_B T \ln g(r)$ from
        PMF calculator (\file{analysis/pmf.hpp})
  \item \textbf{Structural fingerprints}: Topology and geometry hashes for
        classification (\file{classify/fingerprints.cpp})
\end{itemize}

\subsection{Scaling Limits and Atomistic Breakdown}

The atomistic kernel encounters practical limits at system sizes that
motivate the coarse-grained extension:

\begin{center}
\begin{tabular}{lcl}
\toprule
\textbf{Constraint} & \textbf{Threshold} & \textbf{Consequence} \\
\midrule
All-pairs force eval & $N \sim 10{,}000$ & $O(N^2)$ becomes prohibitive \\
Cell-list overhead & $N \sim 100{,}000$ & List management dominates \\
Timestep (H--X bonds) & $\Delta t \leq 0.5$ fs & Limits accessible timescale \\
Equilibration time & $\sim 10^4$ steps & May miss slow relaxation \\
Conformational sampling & $N > 500$ & Barrier-crossing rate drops \\
\bottomrule
\end{tabular}
\end{center}

Beyond these limits, either periodic replication (supercells) or
coarse-graining becomes necessary.  The transition is defined by where
dominant error sources shift, not by arbitrary atom counts (\S10).

\paragraph{Formalized breakdown criterion.}
The atomistic description is no longer efficient when:
\begin{equation}
  \frac{\text{degrees of freedom relevant to target observable}}
       {\text{total atomistic degrees of freedom}} \ll 1
\end{equation}

For bulk diffusion, only center-of-mass motion matters --- individual
bond vibrations are noise.  For phase morphology, only bead-level density
fluctuations matter.  The CG layer systematically eliminates the
irrelevant degrees of freedom.

% =============================================================================
\section{Coarse-Grained Model Construction}
\label{sec:cg_model}
% =============================================================================

This section defines the first real expansion beyond the atomistic layer.
Everything that follows depends on whether this mapping is done with
discipline.

\subsection{Goal}

Convert atomistic outputs into a reduced representation that preserves the
physics relevant to the target observable, while discarding sub-bead
fluctuations.

\subsection{Integration with the Atomistic Layer}

The existing atomistic system becomes the \textbf{reference generator}.
The data flow is:

\begin{enumerate}
  \item Atomistic module generates equilibrated structures / trajectories
  \item Mapping module converts those to bead coordinates
  \item Coarse-graining module estimates reduced interactions
  \item Reduced system is stored in a new but compatible model format
\end{enumerate}

\subsection{Architectural Modules}

Three new module directories are introduced under \file{atomistic/}:

\subsubsection{Module: \texttt{mapping/}}

Responsible for converting atomistic indices to bead indices.

\begin{lstlisting}
// atomistic/mapping/bead_map.hpp
namespace atomistic::mapping {

struct BeadDefinition {
    uint32_t bead_id;
    std::vector<uint32_t> atom_indices;  // atoms in this bead
    std::vector<double> weights;          // mapping weights (mass or uniform)
    std::string label;                    // e.g., "backbone", "sidechain"
};

struct BeadMap {
    std::vector<BeadDefinition> beads;
    uint32_t n_atoms;       // total atomistic atoms
    uint32_t n_beads;       // total CG beads

    // Validate: every atom appears in exactly one bead
    bool is_complete() const;
    bool is_disjoint() const;
};

// Construct bead map from atom-to-bead assignment vector
BeadMap create_bead_map(const State& atomistic_state,
                        const std::vector<uint32_t>& assignments);

} // namespace atomistic::mapping
\end{lstlisting}

\paragraph{Mass-weighted centroid logic.}

The CG bead position $\mathbf{R}_I$ is computed from the atomistic
coordinates via a mapping operator:
\begin{equation}
  \mathbf{R}_I = \sum_{i \in \mathcal{G}_I} w_i \, \mathbf{r}_i
  \label{eq:cg_mapping}
\end{equation}

where $\mathcal{G}_I$ is the set of atoms assigned to bead $I$ and the
weights satisfy $\sum_{i \in \mathcal{G}_I} w_i = 1$.

Two weight schemes are supported:

\begin{center}
\begin{tabular}{ll}
\toprule
\textbf{Scheme} & \textbf{Weights} \\
\midrule
Uniform & $w_i = 1/|\mathcal{G}_I|$ \\
Mass-weighted & $w_i = m_i / \sum_{j \in \mathcal{G}_I} m_j$ \\
\bottomrule
\end{tabular}
\end{center}

Mass-weighted centroids are preferred because they produce bead positions
that coincide with the center of mass of the mapped fragment, preserving
translational dynamics.

\paragraph{Topology compression.}

The atomistic bond graph $\mathbf{B}$ is compressed into a CG connectivity
graph $\mathbf{B}_\text{CG}$:
\begin{equation}
  (I, J) \in \mathbf{B}_\text{CG} \iff \exists\, (i, j) \in \mathbf{B}
  \text{ with } i \in \mathcal{G}_I,\; j \in \mathcal{G}_J,\; I \neq J
\end{equation}

A CG bond exists between beads $I$ and $J$ if any atom in bead $I$ is
bonded to any atom in bead $J$.

\subsubsection{Module: \texttt{cg\_model/}}

Defines the coarse-grained state container and effective interaction terms.

\begin{lstlisting}
// atomistic/cg_model/cg_state.hpp
namespace atomistic::cg_model {

struct CGState {
    uint32_t N_beads;
    std::vector<Vec3> R;         // bead positions (Angstrom)
    std::vector<Vec3> V;         // bead velocities (Ang/fs)
    std::vector<double> M;       // bead masses (amu)
    std::vector<uint32_t> type;  // bead type id

    std::vector<Edge> B_cg;      // CG bond graph
    std::vector<Vec3> F;         // forces on beads
    EnergyTerms E{};             // energy ledger

    BoxPBC box;                  // same PBC as atomistic
};

} // namespace atomistic::cg_model
\end{lstlisting}

\paragraph{Effective pair potentials.}

CG beads interact through effective potentials derived from atomistic
reference data.  The CG pair potential between bead types $A$ and $B$ is:
\begin{equation}
  U_\text{CG}^{AB}(R) = -k_B T \ln g_\text{CG}^{AB}(R) + C
  \label{eq:cg_pair}
\end{equation}

where $g_\text{CG}^{AB}(R)$ is the radial distribution function of
CG beads computed from mapped atomistic trajectories, and $C$ is a
constant chosen so that $U_\text{CG} \to 0$ as $R \to \infty$.

This is the \textbf{potential of mean force} (PMF) --- already implemented
in \file{atomistic/analysis/pmf.hpp} for atomistic pairs and now extended
to CG bead pairs.

\paragraph{Effective bonded terms.}

For bonded CG beads, harmonic approximations are fitted from the
distributions of mapped distances, angles, and dihedrals:
\begin{align}
  U_\text{bond}^\text{CG}(R) &= \tfrac{1}{2} k_R (R - R_0)^2 \\
  U_\text{angle}^\text{CG}(\theta) &= \tfrac{1}{2} k_\theta (\theta - \theta_0)^2
\end{align}

where $R_0$, $\theta_0$ are the equilibrium values and $k_R$, $k_\theta$
are force constants, both extracted from Boltzmann inversion:
\begin{equation}
  k_R = \frac{k_B T}{\langle (R - R_0)^2 \rangle}
  \label{eq:boltzmann_inversion}
\end{equation}

\subsubsection{Module: \texttt{cg\_fit/}}

Provides parameter fitting routines that calibrate CG interactions against
atomistic reference data.

\paragraph{Force matching.}

Given atomistic forces $\mathbf{f}_i$ and the mapping operator, the
instantaneous CG force on bead $I$ is:
\begin{equation}
  \mathbf{F}_I^\text{ref} = \sum_{i \in \mathcal{G}_I} \mathbf{f}_i
  \label{eq:force_matching}
\end{equation}

The CG model parameters $\boldsymbol{\theta}$ are optimized by minimizing:
\begin{equation}
  \chi^2(\boldsymbol{\theta}) =
    \frac{1}{N_\text{frames}} \sum_{t} \sum_{I}
    \left| \mathbf{F}_I^\text{CG}(\boldsymbol{\theta}; \mathbf{R}(t))
         - \mathbf{F}_I^\text{ref}(t) \right|^2
  \label{eq:force_matching_chi2}
\end{equation}

\paragraph{Relative entropy minimization.}

An alternative to force matching.  Minimize the Kullback-Leibler divergence
between the atomistic and CG configuration-space distributions:
\begin{equation}
  S_\text{rel}[\boldsymbol{\theta}]
    = \left\langle \ln \frac{P_\text{AA}(\mathbf{R})}
                             {P_\text{CG}(\mathbf{R}; \boldsymbol{\theta})}
      \right\rangle_\text{AA}
  \label{eq:relative_entropy}
\end{equation}

This is computationally more expensive than force matching but produces
CG models that better reproduce structural distributions.

\paragraph{PMF extraction.}

Already supported via \file{atomistic/analysis/pmf.hpp}.  For CG fitting,
the PMF calculator is applied to mapped bead-pair distances rather than
atomistic pair distances.

\subsection{CG State Compatibility with Engine Conventions}

The \type{CGState} structure mirrors the atomistic \type{State}:

\begin{center}
\begin{tabular}{lll}
\toprule
\textbf{Atomistic} & \textbf{CG Equivalent} & \textbf{Relation} \\
\midrule
\field{State::N} & \field{CGState::N\_beads} & $N_\text{beads} \ll N_\text{atoms}$ \\
\field{State::X} & \field{CGState::R} & Mapped via Eq.~\ref{eq:cg_mapping} \\
\field{State::V} & \field{CGState::V} & Mass-weighted average \\
\field{State::M} & \field{CGState::M} & $M_I = \sum_{i \in \mathcal{G}_I} m_i$ \\
\field{State::type} & \field{CGState::type} & CG bead type ID \\
\field{State::B} & \field{CGState::B\_cg} & Compressed topology \\
\field{State::F} & \field{CGState::F} & CG forces \\
\field{State::E} & \field{CGState::E} & CG energy ledger \\
\field{State::box} & \field{CGState::box} & Same PBC \\
\bottomrule
\end{tabular}
\end{center}

This structural parallel ensures that existing integrators (FIRE, Verlet,
Langevin) can operate on CG systems by adapting the \type{IModel} interface
to CG force evaluation.

\subsection{First Implementation Milestone}

Do not begin with arbitrary chemistry.  Begin with one mapping family:

\begin{enumerate}
  \item Rigid fragment $\to$ one bead (e.g., benzene ring $\to$ one CG bead)
  \item Ligand group $\to$ one bead (e.g., methyl group $\to$ one CH$_3$ bead)
  \item Short repeat unit $\to$ one bead (e.g., ethylene monomer $\to$ one PE bead)
\end{enumerate}

\subsection{Exit Condition}

Section~\ref{sec:cg_model} is integrated when:
\begin{enumerate}
  \item One atomistic system can be mapped to a bead system
  \item The bead system reproduces at least one structural observable
        (e.g., RDF peak position) to within 10\% of atomistic reference
  \item The CG data structure can be simulated using the same high-level
        engine conventions (\type{IModel::eval}, integrators) as atomistic
\end{enumerate}

% =============================================================================
\section{Statistical Mechanics Layer}
\label{sec:stat_mech}
% =============================================================================

This section defines where the project stops returning single ``answers''
and starts returning \textbf{distributions}.

\subsection{Goal}

Wrap both atomistic and CG systems in ensemble-aware execution so that
observables are reported as statistical averages with uncertainties, not
single-point values.

\subsection{Integration with the Current System}

The current engine execution model is:
\begin{enumerate}
  \item Input molecule
  \item Generate candidate
  \item Optimize (FIRE)
  \item Report result
\end{enumerate}

This must expand to:
\begin{enumerate}
  \item Initialize state
  \item Evolve state (MD or MC)
  \item Discard equilibration region
  \item Accumulate observables over production window
  \item Report averages, fluctuations, and distributions
\end{enumerate}

\paragraph{Dual execution modes.}
The existing mode (structure prediction / local minimization) is retained.
A new mode (statistical sampling / time evolution) is added alongside it.
Both modes are selectable through the same high-level run interface.

\subsection{Architectural Modules}

\subsubsection{Module: \texttt{ensembles/}}

Defines ensemble configuration objects.

\begin{lstlisting}
// atomistic/ensembles/ensemble.hpp
namespace atomistic::ensembles {

enum class EnsembleType { NVE, NVT, NPT };

struct EnsembleConfig {
    EnsembleType type = EnsembleType::NVT;
    double temperature = 300.0;    // K (NVT, NPT)
    double pressure = 1.0;         // atm (NPT only)
    uint64_t seed = 42;            // RNG seed for reproducibility
};

} // namespace atomistic::ensembles
\end{lstlisting}

\paragraph{NVE} (microcanonical): Already implemented via velocity Verlet.
Conserves total energy $E = K + U$.

\paragraph{NVT} (canonical): Already implemented via Langevin dynamics.
Maintains target temperature $T$ through friction and stochastic forces.

\paragraph{NPT} (isothermal--isobaric): Not yet implemented.  Requires
a barostat (Berendsen or Parrinello--Rahman) and the pressure conversion
constant $1\,\text{kcal/(mol\,\AA}^3) = 68{,}568.4\,\text{atm}$ (\S4).

\subsubsection{Module: \texttt{samplers/}}

Provides trajectory sampling modes.

\begin{lstlisting}
// atomistic/samplers/md_sampler.hpp
namespace atomistic::samplers {

struct SamplerConfig {
    int n_equilibration = 5000;   // burn-in steps
    int n_production = 50000;     // production steps
    int sample_interval = 100;    // frames between samples
    int checkpoint_interval = 10000;
};

struct SampleResult {
    std::vector<State> frames;         // sampled states
    std::vector<double> energies;      // total energy per frame
    std::vector<double> temperatures;  // instantaneous T per frame
    int n_samples;
    bool equilibrated;                 // stationarity gate passed
};

// MD-based sampler (wraps existing integrators)
SampleResult run_md_sampling(
    State& initial,
    const IModel& model,
    const ModelParams& params,
    const EnsembleConfig& ensemble,
    const SamplerConfig& config
);

} // namespace atomistic::samplers
\end{lstlisting}

A Monte Carlo sampler (\texttt{mc\_sampler.hpp}) is deferred to a later
milestone.

\subsubsection{Module: \texttt{observables/}}

Accumulates running statistics over production trajectories.

\begin{lstlisting}
// atomistic/observables/accumulator.hpp
namespace atomistic::observables {

struct ObservableSet {
    OnlineStats energy_total;
    OnlineStats energy_kinetic;
    OnlineStats temperature;
    OnlineStats pressure;              // when implemented
    OnlineVec3Stats center_of_mass;

    // Derived quantities
    double heat_capacity_Cv() const;   // from energy fluctuations
    double compressibility() const;    // from volume fluctuations (NPT)
};

// Histogram accumulator for distribution observables
class Histogram {
public:
    Histogram(double lo, double hi, int n_bins);
    void add_sample(double x);
    std::vector<double> get_normalized() const;
    double get_mean() const;
    double get_variance() const;
private:
    std::vector<int> counts;
    double lo, hi, bin_width;
    int total;
};

} // namespace atomistic::observables
\end{lstlisting}

\paragraph{Heat capacity from fluctuations.}
In the canonical ensemble:
\begin{equation}
  C_V = \frac{\langle E^2 \rangle - \langle E \rangle^2}{k_B T^2}
  = \frac{\text{Var}(E)}{k_B T^2}
  \label{eq:Cv_fluctuation}
\end{equation}

This requires Welford's algorithm (already in \file{core/statistics.hpp})
to compute $\text{Var}(E)$ in a numerically stable manner.

\subsubsection{Module: \texttt{run\_control/}}

Manages equilibration vs.\ production, checkpointing, and seeded
reproducibility.

\begin{lstlisting}
// atomistic/run_control/run_manager.hpp
namespace atomistic::run_control {

struct RunConfig {
    EnsembleConfig ensemble;
    SamplerConfig sampler;
    ModelParams model_params;

    // Convergence
    int stationarity_window = 10;      // consecutive passes
    double cv_threshold = 0.01;        // coefficient of variation

    // Checkpointing
    std::string checkpoint_dir;
    bool resume_from_checkpoint = false;
};

struct RunResult {
    ObservableSet observables;
    SampleResult samples;
    bool converged;
    int total_steps;
    double wall_time_seconds;
};

RunResult execute(State& initial,
                  const IModel& model,
                  const RunConfig& config);

} // namespace atomistic::run_control
\end{lstlisting}

\subsection{Exit Condition}

Section~\ref{sec:stat_mech} is integrated when:
\begin{enumerate}
  \item The same system can be run as a statistical ensemble rather than
        a one-shot geometry task
  \item Observables are accumulated over many states with stable mean and
        variance
  \item Both atomistic and CG systems plug into the same sampling framework
\end{enumerate}

% =============================================================================
\section{Emergent Behavior Analysis}
\label{sec:emergent}
% =============================================================================

This section sits \textbf{after} the statistical mechanics layer, not
beside it.  Its input is not ``a molecule'' but a \textbf{sample history}
or \textbf{ensemble dataset}.

\subsection{Goal}

Extract physically meaningful collective behavior from ensemble output:
bulk structure, not individual coordinates.

\subsection{Data Flow}

\begin{enumerate}
  \item Section~\ref{sec:stat_mech} produces trajectories / sampled states
  \item This section reads those states
  \item Computes collective observables
  \item Report layer exports \textit{trends}, not just snapshots
\end{enumerate}

\subsection{Architectural Modules}

\subsubsection{Module: \texttt{trajectory/}}

Frame storage, state readers, and decimation.

\begin{lstlisting}
// atomistic/trajectory/trajectory.hpp
namespace atomistic::trajectory {

class Trajectory {
public:
    void add_frame(const State& s, double time);
    const State& frame(int i) const;
    int n_frames() const;

    // Decimation: keep every k-th frame
    Trajectory decimate(int k) const;

    // Time window: frames between t_start and t_end
    Trajectory window(double t_start, double t_end) const;

private:
    std::vector<State> frames_;
    std::vector<double> times_;
};

} // namespace atomistic::trajectory
\end{lstlisting}

\subsubsection{Module: \texttt{analysis/structure/}}

Structural collective observables.

\paragraph{Radial distribution function $g(r)$.}
Already partially implemented in \file{analysis/pmf.cpp} for PMF
extraction.  The standalone RDF calculation is:
\begin{equation}
  g(r) = \frac{V}{N^2} \left\langle \sum_{i \neq j}
         \delta(r - r_{ij}) \right\rangle \bigg/ (4\pi r^2 \Delta r)
  \label{eq:rdf}
\end{equation}

\paragraph{Nearest-neighbor statistics.}
Distribution of coordination numbers $n_i = |\{j : r_{ij} < r_\text{cut}\}|$
for each atom type.

\paragraph{Clustering.}
Already implemented in \file{classify/cluster.cpp}.  For trajectory analysis,
cluster size distributions are computed per frame and averaged:
\begin{equation}
  \langle n_s \rangle = \frac{1}{N_\text{frames}}
  \sum_{t=1}^{N_\text{frames}} n_s(t)
  \label{eq:cluster_avg}
\end{equation}

where $n_s(t)$ is the number of clusters of size $s$ at frame $t$.

\subsubsection{Module: \texttt{analysis/order/}}

Order parameters and density fluctuations.

\paragraph{Bond-orientational order parameter $Q_l$.}
For characterizing local crystalline order:
\begin{equation}
  Q_l(i) = \left( \frac{4\pi}{2l+1}
            \sum_{m=-l}^{l} |q_{lm}(i)|^2 \right)^{1/2}
  \label{eq:ql}
\end{equation}

where $q_{lm}(i) = \frac{1}{n_i} \sum_{j=1}^{n_i} Y_{lm}(\hat{\mathbf{r}}_{ij})$
and $Y_{lm}$ are spherical harmonics.  $Q_6$ distinguishes FCC, BCC, HCP,
and liquid phases.

\paragraph{Density fluctuations.}
Local number density in subvolumes $\Omega_k$:
\begin{equation}
  \rho_k = \frac{N_k}{|\Omega_k|}
\end{equation}

The variance $\text{Var}(\rho)$ over subvolumes measures structural
heterogeneity.

\subsubsection{Module: \texttt{analysis/free\_energy/}}

Histogram projections and PMF over reaction coordinates.

\begin{equation}
  F(\xi) = -k_B T \ln P(\xi) + C
  \label{eq:free_energy_1d}
\end{equation}

where $\xi$ is a collective variable (e.g., cluster size, density, distance)
and $P(\xi)$ is its probability distribution accumulated over the production
trajectory.

\subsection{First Implementation Milestone}

Implement only:
\begin{enumerate}
  \item RDF $g(r)$ for arbitrary pair types
  \item Cluster size distribution
  \item One order parameter ($Q_6$)
\end{enumerate}

This alone gives real bulk-property language.

\subsection{Exit Condition}

Section~\ref{sec:emergent} is integrated when:
\begin{enumerate}
  \item The system can produce ensemble trajectories
  \item Those trajectories can be reduced into structural observables
  \item Those observables change meaningfully with temperature, density,
        or interaction strength
\end{enumerate}

% =============================================================================
\section{Multi-Scale Coupling}
\label{sec:multiscale}
% =============================================================================

This section formalizes how information passes consistently between
atomistic, coarse-grained, and continuum-like descriptions.

\subsection{Goal}

Allow verified, repeatable inter-scale information transfer.  This means
explicit operators --- not vague claims of ``multi-scale capability.''

\subsection{Sequential Coupling (First Target)}

\begin{enumerate}
  \item Atomistic run produces local physics (energies, forces, $g(r)$)
  \item Fit / infer CG properties (PMF, effective potentials)
  \item Run CG simulation at extended scale
  \item Extract effective bulk coefficients (density, diffusion, elastic moduli)
  \item Pass to continuum description (if applicable)
\end{enumerate}

\textbf{Concurrent coupling} (atomistic subregion embedded in CG or
continuum, adaptive handoff during runtime) is deferred.  It is a
qualitatively harder problem requiring force interpolation at boundaries,
energy conservation across scales, and a thermostat that maintains correct
temperature in both regions.

\subsection{Architectural Modules}

\subsubsection{Module: \texttt{multiscale/bridges/}}

\paragraph{Atomistic $\to$ CG bridge.}

\begin{lstlisting}
// atomistic/multiscale/bridges/aa_to_cg.hpp
namespace atomistic::multiscale {

struct CGBridgeResult {
    cg_model::CGState cg_state;
    mapping::BeadMap bead_map;
    std::vector<double> effective_pair_potential;
    std::vector<double> effective_bond_potential;
};

// Map a single atomistic state to CG
CGBridgeResult map_to_cg(
    const State& atomistic_state,
    const mapping::BeadMap& map
);

// Fit CG potentials from atomistic trajectory
CGBridgeResult fit_cg_model(
    const trajectory::Trajectory& aa_trajectory,
    const mapping::BeadMap& map,
    double temperature
);

} // namespace atomistic::multiscale
\end{lstlisting}

\paragraph{CG $\to$ continuum bridge.}

Extracts effective scalar and tensor properties from CG simulation:

\begin{itemize}
  \item Average density $\bar{\rho} = N M_\text{bead} / V$
  \item Diffusion coefficient $D$ from mean-squared displacement:
        $D = \lim_{t \to \infty} \frac{\langle |\mathbf{R}(t) - \mathbf{R}(0)|^2 \rangle}{6t}$
  \item Elastic moduli from stress fluctuations (Born--Huang method)
\end{itemize}

\paragraph{Backmapping (optional, deferred).}

Reconstruct atomistic detail from CG configuration for local refinement.
This requires fragment libraries and constrained energy minimization.

\subsubsection{Module: \texttt{homogenization/}}

\begin{lstlisting}
// atomistic/homogenization/bulk_properties.hpp
namespace atomistic::homogenization {

struct BulkProperties {
    double density;                    // g/cm^3
    double diffusion_coefficient;      // cm^2/s
    double bulk_modulus;               // GPa (from volume fluctuations)
    double shear_modulus;              // GPa (from stress tensor)
    double thermal_conductivity;       // W/(m*K) (from Green-Kubo)
};

BulkProperties extract_from_trajectory(
    const trajectory::Trajectory& traj,
    double temperature,
    double volume
);

} // namespace atomistic::homogenization
\end{lstlisting}

\subsubsection{Module: \texttt{field\_export/}}

Projects particle data onto continuum-compatible grid representations.

\begin{itemize}
  \item Local density maps on regular grids
  \item Coarse field tensors (stress, velocity gradient)
  \item Suitable for coupling with finite element or finite difference solvers
\end{itemize}

\subsection{First Implementation Milestone}

Implement only:
\begin{enumerate}
  \item Atomistic $\to$ CG mapping and potential fitting
  \item CG $\to$ effective scalar observables (density, diffusion)
\end{enumerate}

No runtime hybrid coupling.

\subsection{Exit Condition}

Section~\ref{sec:multiscale} is integrated when:
\begin{enumerate}
  \item Properties inferred from smaller-scale models can parameterize
        larger-scale runs
  \item The data pathway is explicit and repeatable
  \item One property can be shown propagating correctly across scales
\end{enumerate}

% =============================================================================
\section{Validation and Benchmarks}
\label{sec:validation}
% =============================================================================

This section co-evolves with Sections~\ref{sec:cg_model} through
\ref{sec:multiscale}.  It is not an afterthought.

\subsection{Goal}

Continuously verify that each added scale still reproduces relevant physics.

\subsection{Benchmark Suite}

\begin{center}
\begin{tabular}{lcl}
\toprule
\textbf{Benchmark} & \textbf{Source Section} & \textbf{Observable} \\
\midrule
LJ Argon (atomistic) & \S3 & Energy, $g(r)$, $T$ \\
LJ Argon (mapped CG) & \ref{sec:cg_model} & CG $g(r)$ vs.\ atomistic \\
NVT equilibration & \ref{sec:stat_mech} & $\langle T \rangle$, $\text{Var}(E)$ \\
RDF convergence & \ref{sec:emergent} & $g(r)$ peak position vs.\ $N_\text{frames}$ \\
Scale-transfer fidelity & \ref{sec:multiscale} & CG density vs.\ atomistic \\
\bottomrule
\end{tabular}
\end{center}

\subsection{Metrics}

\begin{itemize}
  \item \textbf{RMSE}: Root mean square error between CG and atomistic $g(r)$
  \item \textbf{Relative error}: $|x_\text{CG} - x_\text{AA}| / |x_\text{AA}|$
        for scalar observables
  \item \textbf{Structural similarity}: Overlap integral of RDFs
  \item \textbf{Diffusion coefficient error}: Relative error in $D$ between
        CG and atomistic
  \item \textbf{Finite-size deviation}: Observable shift with system size
\end{itemize}

\subsection{Regression Testing}

Seeded standard runs produce deterministic outputs.  A regression test
compares current output against stored reference:

\begin{lstlisting}
// tests/cg_regression_test.cpp
TEST(CGRegression, ArgonCG_RDF) {
    auto aa_traj = load_reference_trajectory("argon_108_nvt.xyzf");
    auto bead_map = create_single_atom_map(aa_traj.frame(0));
    auto cg_result = fit_cg_model(aa_traj, bead_map, 85.0);

    auto rdf = compute_rdf(cg_result.cg_state, 0.1, 10.0);
    auto ref = load_reference_rdf("argon_cg_rdf_ref.csv");

    double rmse = compute_rmse(rdf, ref);
    EXPECT_LT(rmse, 0.05);  // 5% tolerance
}
\end{lstlisting}

\subsection{Exit Condition}

Section~\ref{sec:validation} is integrated when:
\begin{enumerate}
  \item Every phase from \ref{sec:cg_model} onward has at least one benchmark
  \item Changes in code can be checked automatically against expected
        physical outputs
  \item Loss of fidelity can be quantified rather than hand-waved
\end{enumerate}

% =============================================================================
\section{Scale Transition: Beyond $10^{-9}$\,m}
\label{sec:scale_transition}
% =============================================================================

This section formalizes when the representation itself must change.

\subsection{Goal}

Define operational thresholds for leaving atomistic detail behind.

\subsection{Scale-Aware Control Logic}

\begin{lstlisting}
// atomistic/scale_manager/regime_selector.hpp
namespace atomistic::scale_manager {

enum class Regime {
    Atomistic,          // Full atomistic detail
    CoarseGrained,      // Reduced CG representation
    MultiscaleSeq,      // Sequential atomistic -> CG -> properties
    ContinuumExport     // Export effective properties only
};

struct ProblemSpec {
    int n_atoms;
    double domain_length;       // Angstrom
    double target_timescale;    // fs
    std::string target_observable;
};

Regime recommend_regime(const ProblemSpec& spec);

} // namespace atomistic::scale_manager
\end{lstlisting}

\paragraph{Rules-based selector (first implementation).}

\begin{center}
\begin{tabular}{lcl}
\toprule
\textbf{Condition} & \textbf{Recommended Regime} \\
\midrule
Discrete compact molecule, $N < 500$ & Atomistic \\
Many repeated units, $500 < N < 10{,}000$ & CG \\
$N > 10{,}000$ or domain $> 100$\,\AA & MultiscaleSeq \\
Domain $> 1{,}000$\,\AA & ContinuumExport \\
\bottomrule
\end{tabular}
\end{center}

\subsection{Exit Condition}

Section~\ref{sec:scale_transition} is integrated when:
\begin{enumerate}
  \item Given a physical problem, the code can recommend an appropriate
        representation
  \item Scale transition rules are explicit and tied to cost and physics
  \item The thesis can defend why different regimes use different models
\end{enumerate}

% =============================================================================
\section{Thermostats and Stochastic Dynamics}
\label{sec:thermostats}
% =============================================================================

This section upgrades the integrator interface so that force evaluation is
no longer the only contributor to state evolution.

\subsection{Goal}

Introduce realistic environmental exchange and fluctuation physics.

\subsection{Current State}

Langevin dynamics is already implemented in
\file{atomistic/integrators/velocity\_verlet.cpp} with:
\begin{itemize}
  \item Friction coefficient $\gamma$
  \item Random force $\sqrt{2\gamma m_i k_B T}\;\xi_i$ where $\xi_i \sim \mathcal{N}(0,1)$
  \item Temperature error $< 0.6\%$ validated
\end{itemize}

\subsection{Enhancements}

\subsubsection{Thermostat abstraction layer}

\begin{lstlisting}
// atomistic/thermostats/thermostat.hpp
namespace atomistic::thermostats {

struct IThermostat {
    virtual ~IThermostat() = default;
    virtual void apply(State& s, double dt) = 0;
    virtual double target_temperature() const = 0;
};

// Velocity rescale (simple, for testing only)
class VelocityRescale : public IThermostat { ... };

// Langevin (primary, already implemented)
class Langevin : public IThermostat { ... };

// Nose-Hoover (extended Lagrangian, deferred)
class NoseHoover : public IThermostat { ... };

} // namespace atomistic::thermostats
\end{lstlisting}

\subsubsection{Noise generation}

\begin{lstlisting}
// atomistic/noise/gaussian_rng.hpp
namespace atomistic::noise {

class GaussianRNG {
public:
    explicit GaussianRNG(uint64_t seed);
    double sample();                    // N(0,1)
    Vec3 sample_vec3();                 // 3 independent N(0,1)
    void validate_covariance(int n);    // statistical self-test
private:
    std::mt19937_64 gen_;
    std::normal_distribution<double> dist_;
};

} // namespace atomistic::noise
\end{lstlisting}

\subsection{Timestep Logic}

The enhanced integrator timestep becomes:
\begin{enumerate}
  \item Deterministic forces: $\mathbf{F}_i = -\nabla_i U$
  \item Thermostat correction: friction $-\gamma m_i \mathbf{v}_i$
  \item Stochastic term: $\sqrt{2\gamma m_i k_B T / \Delta t}\;\boldsymbol{\xi}_i$
  \item State update: position and velocity advance
\end{enumerate}

\subsection{Exit Condition}

Section~\ref{sec:thermostats} is integrated when:
\begin{enumerate}
  \item Target temperature is maintained statistically ($< 1\%$ error)
  \item Fluctuations are physically consistent (equipartition satisfied)
  \item Deterministic and stochastic runs can be selected through the
        same high-level run interface
\end{enumerate}

% =============================================================================
\section{Free Energy and Phase Behavior}
\label{sec:free_energy}
% =============================================================================

This is the top analytical layer.  It consumes the work of
Sections~\ref{sec:stat_mech}, \ref{sec:emergent}, and \ref{sec:thermostats}.

\subsection{Goal}

Move from ``what structure appeared'' to ``what states are thermodynamically
preferred.''

\subsection{Collective Variables}

\begin{lstlisting}
// atomistic/collective_variables/cv.hpp
namespace atomistic::collective_variables {

struct ICollectiveVariable {
    virtual ~ICollectiveVariable() = default;
    virtual double compute(const State& s) const = 0;
    virtual std::string name() const = 0;
};

// Concrete CVs
class NumberDensity : public ICollectiveVariable { ... };
class ClusterFraction : public ICollectiveVariable { ... };
class CoordinationNumber : public ICollectiveVariable { ... };
class CompositionFraction : public ICollectiveVariable { ... };

} // namespace atomistic::collective_variables
\end{lstlisting}

\subsection{Free Energy Reconstruction}

\paragraph{Histogram method.}

Accumulate $P(\xi)$ over the production trajectory:
\begin{equation}
  F(\xi) = -k_B T \ln P(\xi) + C
\end{equation}

This gives the free energy as a function of collective variable $\xi$.
The constant $C$ is chosen so that $\min_\xi F(\xi) = 0$.

\paragraph{PMF reconstruction.}

For pairwise collective variables (e.g., bead--bead distance), the existing
PMF calculator (\file{analysis/pmf.hpp}) is reused directly.

\paragraph{Reweighting (deferred).}

Weighted histogram analysis method (WHAM) for combining data from
simulations at different thermodynamic conditions.

\subsection{Phase Analysis}

\begin{lstlisting}
// atomistic/phase_analysis/phase.hpp
namespace atomistic::phase_analysis {

struct PhaseResult {
    std::vector<double> cv_values;
    std::vector<double> free_energy;
    std::vector<int> phase_labels;      // per-frame phase assignment
    int n_phases_detected;
    double transition_temperature;       // if applicable
};

// Detect phase transitions from free energy landscape
PhaseResult analyze_phases(
    const trajectory::Trajectory& traj,
    const ICollectiveVariable& cv,
    double temperature,
    int n_bins = 100
);

} // namespace atomistic::phase_analysis
\end{lstlisting}

\subsection{First Implementation Milestone}

Do not begin with full phase diagrams.  Begin with:
\begin{enumerate}
  \item One order parameter (e.g., $Q_6$)
  \item One histogram accumulation
  \item One PMF or free-energy projection
  \item Two temperatures or densities for comparison
\end{enumerate}

\subsection{Exit Condition}

Section~\ref{sec:free_energy} is integrated when:
\begin{enumerate}
  \item Sampled state populations can be converted into free-energy-like
        surfaces
  \item Changes in thermodynamic conditions visibly shift state preference
  \item Phase-like transitions can be described from simulation outputs
        rather than intuition
\end{enumerate}

% =============================================================================
\section{Summary: Dependency Graph and Implementation Order}
% =============================================================================

The sections form a strict dependency chain:

\begin{center}
\begin{tabular}{clcl}
\toprule
\textbf{Order} & \textbf{Section} & \textbf{Depends On} & \textbf{Key Deliverable} \\
\midrule
1 & \ref{sec:cg_model} CG Model & Atomistic (existing) & Bead mapping + PMF fitting \\
2 & \ref{sec:stat_mech} Stat.\ Mech. & Atomistic + CG & Ensemble sampling framework \\
3 & \ref{sec:emergent} Emergent & \ref{sec:stat_mech} & $g(r)$, clusters, $Q_6$ \\
4 & \ref{sec:multiscale} Multi-Scale & \ref{sec:cg_model} + \ref{sec:emergent} & AA$\to$CG$\to$bulk bridge \\
5 & \ref{sec:validation} Validation & All above & Benchmark suite \\
6 & \ref{sec:scale_transition} Scale Trans. & \ref{sec:multiscale} & Regime recommender \\
7 & \ref{sec:thermostats} Thermostats & \ref{sec:stat_mech} & Thermostat abstraction \\
8 & \ref{sec:free_energy} Free Energy & \ref{sec:emergent} + \ref{sec:thermostats} & Phase analysis \\
\bottomrule
\end{tabular}
\end{center}

\paragraph{Critical path.}
Sections \ref{sec:cg_model} $\to$ \ref{sec:stat_mech} $\to$
\ref{sec:emergent} form the critical path.  Everything else can proceed in
parallel once Section~\ref{sec:stat_mech} is operational.

\paragraph{Validation is continuous.}
Section~\ref{sec:validation} is not a final step.  Every module added must
include at least one benchmark test.  The validation suite grows with the
codebase.

% =============================================================================
\section{Namespace and Directory Plan}
% =============================================================================

All new code resides under the existing \file{atomistic/} directory to
maintain the single-library build model.  New subdirectories:

\begin{center}
\begin{tabular}{ll}
\toprule
\textbf{Directory} & \textbf{Purpose} \\
\midrule
\file{atomistic/mapping/} & Atom-to-bead maps, centroid logic, topology compression \\
\file{atomistic/cg\_model/} & Bead types, masses, connectivity, effective potentials \\
\file{atomistic/cg\_fit/} & Force matching, relative entropy, PMF extraction \\
\file{atomistic/ensembles/} & NVE, NVT, NPT configuration objects \\
\file{atomistic/samplers/} & MD runner, MC runner (future) \\
\file{atomistic/observables/} & Accumulators, histograms, correlation functions \\
\file{atomistic/run\_control/} & Burn-in, production, checkpointing \\
\file{atomistic/trajectory/} & Frame storage, decimation, windowing \\
\file{atomistic/analysis/structure/} & $g(r)$, nearest-neighbor, clustering \\
\file{atomistic/analysis/order/} & Order parameters, density fluctuations \\
\file{atomistic/analysis/free\_energy/} & Histogram, PMF, reweighting \\
\file{atomistic/multiscale/bridges/} & AA$\to$CG, CG$\to$continuum, backmapping \\
\file{atomistic/homogenization/} & Density, diffusion, elastic moduli \\
\file{atomistic/field\_export/} & Grid projection, density maps, field tensors \\
\file{atomistic/thermostats/} & Velocity rescale, Langevin, Nos\'{e}--Hoover \\
\file{atomistic/noise/} & Gaussian RNG, covariance validation \\
\file{atomistic/collective\_variables/} & Density, cluster fraction, coordination \\
\file{atomistic/phase\_analysis/} & Coexistence, transitions, free-energy surfaces \\
\file{atomistic/scale\_manager/} & Problem classification, regime selection \\
\file{atomistic/benchmarks/} & Reference systems, known observables, comparators \\
\file{atomistic/metrics/} & RMSE, relative error, structural similarity \\
\bottomrule
\end{tabular}
\end{center}

All modules link against the existing \type{atomistic} static library
target.  New source files are added to \file{atomistic/CMakeLists.txt}.

% =============================================================================
\section{Compatibility Contract}
% =============================================================================

\subsection{With Existing Atomistic Layer}

\begin{enumerate}
  \item The \type{State} struct is not modified.  CG representations use
        \type{CGState}, a parallel structure.
  \item The \type{IModel} interface is not modified.  CG force models
        implement a \type{ICGModel} interface with the same
        \op{eval(CGState\&, ...)} signature pattern.
  \item Existing integrators (FIRE, Verlet, Langevin) are not modified.
        Template or adapter wrappers allow them to operate on \type{CGState}.
  \item File formats are extended, not replaced.  CG states use
        \texttt{.xyz\_cg} or the reserved CG slots in \texttt{.xyzc}.
  \item The bridge architecture (\S SBA) is preserved.  CG outputs pass
        through the same canonical structural bridge to the desktop layer.
\end{enumerate}

\subsection{With Build System}

\begin{enumerate}
  \item All CG code compiles as part of the \type{atomistic} static library
  \item No new external dependencies
  \item C++20 standard maintained
  \item No fast-math flags (determinism contract, \S11)
\end{enumerate}

\subsection{With Validation Framework}

\begin{enumerate}
  \item Every new module includes at least one unit test
  \item Every CG observable has a regression test against atomistic reference
  \item Seeded runs produce deterministic outputs (same seed $\Rightarrow$
        bit-identical result)
\end{enumerate}

% =============================================================================
\section{Emergent Effective Medium Mapping}
\label{sec:ensemble_proxy}
% =============================================================================

This section documents the \textbf{Ensemble-Level Macroscopic Response Proxy}
layer (\file{coarse\_grain/analysis/ensemble\_proxy.hpp}).  It is the bridge
between per-bead environment state and interpretable bulk characterisation.

Architecture position:
\begin{center}
\begin{tabular}{c}
Environment-state evaluation \\
$\downarrow$ \\
Ensemble statistics / spatial field summaries \\
$\downarrow$ \\
Macroscopic response proxies \\
$\downarrow$ \\
Later constitutive or transport models
\end{tabular}
\end{center}

\subsection{Anti-Black-Box Commitment}

Every proxy is an explicit, inspectable function of bead-state distributions.
No hidden weights.  No learned parameters.  All intermediate statistics are
exposed in \type{EnsembleProxySummary}.  This section records the design
decisions that make the layer traceable.

\subsection{Model Parameters}

Three model-level constants are declared at namespace scope:

\begin{description}
  \item[\texttt{PROXY\_MIN\_BEADS = 8}]
    Minimum ensemble size before proxies are considered reliable.
    Below this, all proxy fields are still computed but the
    \type{EnsembleProxySummary::valid} flag is set to \texttt{false}.
    \textit{Rationale}: variance and autocorrelation estimates are unreliable
    for $N < 8$; bulk/edge classification degenerates at small~$N$.

  \item[\texttt{PROXY\_CONVERGENCE\_THRESH = 1e-4}]
    Threshold for the \fn{compute\_proxy\_delta} convergence flag.
    The system is declared converged when
    $|\Delta\bar\eta| < \epsilon$ \textbf{and}
    $|\Delta\overline{|\eta - f|}| < \epsilon$.

  \item[\texttt{PROXY\_VARIANCE\_FLOOR = 1e-12}]
    Lower bound applied in proxy denominators to prevent division-by-zero
    on degenerate (all-identical) inputs.  Does not affect the stored
    \texttt{var\_*} fields.
\end{description}

These are explicit model parameters, not magic numbers buried in formulas.

\subsection{Proxy Formulas}
\label{sec:proxy_formulas}

All five proxies produce values in $[0, 1]$.  Their explicit definitions are:

\subsubsection{Cohesion proxy}
\begin{equation}
  \text{cohesion} = \frac{\bar\rho_{\hat{}} + \bar\eta
    + (1 - \overline{|\eta - f|})}{3}
  \label{eq:cohesion_proxy}
\end{equation}

Combines density response $\bar\rho_{\hat{}}$, internal stabilisation
$\bar\eta$, and convergence quality $(1 - \text{mean mismatch})$.
Higher $\Rightarrow$ more cohesive (denser, more adapted, more converged).
\textit{Does not} measure cohesive energy.

\subsubsection{Uniformity proxy}
\begin{equation}
  \text{uniformity} = 1 - \frac{1}{2}\left(
    \min(2\sigma_\eta,\, 1) + \min(2\sigma_{\rho_{\hat{}}},\, 1)\right)
  \label{eq:uniformity_proxy}
\end{equation}

where $\sigma = \sqrt{\text{var}}$.  For $[0,1]$-bounded variables the
maximum standard deviation is $\approx 0.5$; multiplying by~2 maps
$[0, 0.5] \to [0, 1]$.  Higher $\Rightarrow$ more homogeneous packing.

\subsubsection{Texture proxy}
\begin{equation}
  \text{texture} = \bar{P}_{2,\hat{}}
  \label{eq:texture_proxy}
\end{equation}

Orientational organisation as measured by the mean normalised second
Legendre polynomial.  Range: $0$ (perpendicular / no alignment),
$\approx 0.33$ (isotropic random), $1$ (perfectly aligned).

\subsubsection{Stabilisation proxy}
\begin{equation}
  \text{stabilisation} = \bar\eta \cdot (1 - \overline{|\eta - f|})
  \label{eq:stabilization_proxy}
\end{equation}

Adaptation quality.  High only when $\bar\eta$ is high \textit{and}
the mean mismatch is low.  High mismatch can indicate either ``still
converging'' or ``actively diverging'' --- distinguishable only via the
time-delta fields (Section~\ref{sec:proxy_time_history}).

\subsubsection{Surface sensitivity proxy}
\begin{equation}
  \text{surface\_sens} = \frac{1}{3}\left(
    \frac{|\Delta_\rho|}{\bar\rho + \varepsilon} +
    \frac{|\Delta_C|}{\bar{C} + \varepsilon} +
    \frac{|\Delta_\eta|}{\bar\eta + \varepsilon} \right)
  \label{eq:surface_sensitivity_proxy}
\end{equation}

where $\Delta_x = x_\text{bulk} - x_\text{edge}$ are the bulk-minus-edge
gaps from the classifier (Section~\ref{sec:bulk_edge_classifier}).
Higher $\Rightarrow$ more surface-dominated structure.

\textit{Guard}: when the gap fields are NaN (no-bulk or no-edge case),
surface sensitivity is set to~$0$.

\subsection{Bulk/Edge Classifier}
\label{sec:bulk_edge_classifier}

The function \fn{classify\_bulk\_edge()} applies two criteria.  \textbf{OR
logic}: failing either criterion marks a bead as edge (surface).

\paragraph{Criterion 1 — Coordination (always active).}
A bead with $C_i < C_\text{thresh}$ is edge.
$C_\text{thresh}$ is a \textbf{model parameter}:
\begin{itemize}
  \item Default (\texttt{coord\_threshold < 0}): auto-median of the
        ensemble $C$ distribution.
  \item Explicit value: the caller pins $C_\text{thresh}$ directly.
        Required for comparative studies.
  \item Traceability: the resolved value is written to
        \texttt{EnsembleProxySummary::c\_thresh\_used}.
\end{itemize}

\paragraph{Criterion 2 — Geometric depth (optional).}
When \texttt{surface\_depth > 0}, any bead within that distance of the
radial boundary of the ensemble centroid is additionally classified as edge,
regardless of coordination.
\textit{Precedence}: OR logic is conservative --- more beads are classified
as surface, not fewer.  This prevents false bulk classification near
boundaries.

\paragraph{Gap field guards.}
When $n_\text{bulk} = 0$ or $n_\text{edge} = 0$ the gap fields
($\Delta_\rho$, $\Delta_C$, $\Delta_{P_2}$, $\Delta_\eta$) are set to
\texttt{NaN}, not zero.  Zero would incorrectly imply ``no gap''; NaN
correctly signals ``undefined''.

\subsection{Spatial Correlation Length}
\label{sec:xi_extraction}

\fn{spatial\_correlation\_length()} estimates $\xi$ by fitting the
empirical pair correlation $C(r)$ to an exponential decay
$C(r) \propto \exp(-r/\xi)$ via log-linear regression.

\paragraph{Method.}
\begin{enumerate}
  \item Bin all pairs by separation $r$ into bins of width $\Delta r = 1$\,\AA.
  \item For each bin with $\geq 3$ pairs, compute
        $C(r_k) = \langle (x_i - \bar x)(x_j - \bar x) \rangle / \sigma^2$.
  \item For bins where $C(r_k) > 0$, fit $\ln C(r)$ against $r$ by
        ordinary linear regression (slope $b = -1/\xi$).
  \item Return $\xi = -1/b$.
\end{enumerate}

\paragraph{Failure modes (returns NaN).}
\begin{itemize}
  \item $\sigma^2 < 10^{-30}$: field is uniform; $\xi$ undefined.
  \item Fewer than 2 bins have enough pairs.
  \item All $C(r_k) \leq 0$: correlation does not decay log-linearly.
  \item Regression slope $b \geq 0$: field is not decaying.
\end{itemize}

\paragraph{Edge participation.}
All beads participate regardless of bulk/edge class.  Excluding edge beads
changes $\xi$ systematically; selective exclusion is the caller's
responsibility via subsetting.

\subsection{Time History and Convergence}
\label{sec:proxy_time_history}

\fn{compute\_proxy\_delta(curr, prev)} fills the delta fields of a
current snapshot given a previous one:
\begin{align}
  \Delta\bar\eta &= \bar\eta^\text{curr} - \bar\eta^\text{prev} \\
  \Delta\overline{|\eta-f|} &= \overline{|\eta-f|}^\text{curr}
    - \overline{|\eta-f|}^\text{prev}
\end{align}

The \texttt{converged} flag is set when both deltas fall below
\texttt{PROXY\_CONVERGENCE\_THRESH}.

\textit{Why this matters}: a high mismatch in a single snapshot could mean
``still converging'' ($\Delta > 0$, improving) or ``actively diverging''
($\Delta < 0$, worsening) --- opposite physical situations requiring
opposite responses.  Without $\Delta$ fields, the proxy is uninterpretable.

\subsection{Reference Distributions and Normalisation}

Three reference configurations should be run before any production analysis:
\begin{enumerate}
  \item \textbf{Ideal lattice} at target $N$ and density: structural upper
        bound for uniformity and cohesion.
  \item \textbf{Randomised positions}, same $N$ and density: lower bound.
  \item \textbf{Structured positions, random orientations}: lower bound for
        texture.
\end{enumerate}

Stored in \type{EnsembleProxyReference\{ideal, random, disordered\}}.

The function \fn{apply\_reference(s, lo, hi)} fills the \texttt{rel\_*}
fields of a target summary:
\begin{equation}
  \text{rel} = \clamp\!\left(
    \frac{v - v_\text{lo}}{v_\text{hi} - v_\text{lo}},\; 0,\; 1\right)
\end{equation}

If $v_\text{hi} \approx v_\text{lo}$ (degenerate reference), \texttt{rel}
is set to $0.5$ (indeterminate).

\paragraph{Why this is necessary.}
A raw proxy value of $0.73$ is meaningless without a baseline.
Reference-normalised values place the system on an anchored $[0, 1]$
scale where $0$ = random reference and $1$ = ideal lattice reference.

\subsection{Numeric Guards}

The following guards are applied:
\begin{description}
  \item[Validity flag]
    \texttt{EnsembleProxySummary::valid = false} for $N < \texttt{PROXY\_MIN\_BEADS}$.
    Proxy fields are still computed so that callers can inspect them,
    but the valid flag signals that quantitative interpretation is unreliable.

  \item[Variance floor]
    Proxy denominators apply \texttt{PROXY\_VARIANCE\_FLOOR} to prevent
    division by zero on all-identical inputs.
    Stored \texttt{var\_*} fields are exact (unmodified).

  \item[Gap field NaN]
    When no bulk beads or no edge beads exist, all four gap fields
    are set to \texttt{NaN}.  The surface sensitivity proxy treats NaN
    gaps as zero sensitivity (the only physically defensible default
    when the classification is degenerate).
\end{description}

\subsection{Directed Proxy Validation (Suite 5F)}

Each proxy has a mandatory directed test that specifies the perturbation,
the expected direction of change, and a strict pass criterion.
If any test fails, the proxy formula is wrong and must be fixed before
use in any result.

\begin{center}
\begin{tabular}{lll}
\textbf{Proxy} & \textbf{Perturbation} & \textbf{Pass criterion} \\
\hline
Cohesion & Degrade edge bead $\rho_\hat{}, \eta$ to zero & cohesion decreases \\
Texture & Set all $\hat{P}_2 = 1/3$ & texture\_proxy $= 1/3$ exactly \\
Uniformity & Set all states to ensemble mean & uniformity $> 0.99$ \\
Stabilisation & Freeze $\eta$ at random, target\_f$=0.9$ & stabilisation decreases \\
Surface sens. & All beads identical & surface sensitivity $< 0.02$ \\
\end{tabular}
\end{center}

All five tests pass as of Suite \#5 (131/131).

% =============================================================================
\section{Macro Property Precursor Layer}
\label{sec:macro_precursor}
% =============================================================================

This section documents the \textbf{Macro Property Precursor Channel}
layer (\file{coarse\_grain/analysis/macro\_precursor.hpp}).
It converts ensemble-level macroscopic response proxies into property
precursor bundles: intermediate latent descriptors grouped by
material-behaviour hypothesis.

\subsection{Why a Precursor Layer Is Needed}

Coarse-grained ensemble statistics do not directly produce physical
engineering properties.  The gap between ``high cohesion'' and
``tensile strength = 430~MPa'' is enormous.  This layer does
\textbf{not} close that gap.  Instead it:

\begin{enumerate}
  \item Scrapes precursor signals from the proxy layer
  \item Assembles them into candidate property channels
  \item Tests directional correctness (monotone response)
  \item Preserves provenance (contributing proxies, weights, validity)
  \item Prepares later calibration against real data
\end{enumerate}

The pipeline is:
\begin{center}
\begin{tabular}{c}
Environment-state evaluation \\
$\downarrow$ \\
Ensemble statistics / spatial field summaries \\
$\downarrow$ \\
Macroscopic response proxies \\
$\downarrow$ \\
Macro precursor channels \quad $\leftarrow$ \textit{this layer} \\
$\downarrow$ \\
Later calibrated property estimator
\end{tabular}
\end{center}

\subsection{Naming Convention}

All channel names carry the suffix \texttt{-like} to indicate that they
are \textbf{evidence variables}, not physical properties:

\begin{center}
\begin{tabular}{ll}
\textbf{Channel} & \textbf{Physical hypothesis} \\
\hline
\texttt{rigidity\_like}           & stiff, well-connected bulk \\
\texttt{ductility\_like}          & deformable without fracture \\
\texttt{brittleness\_like}        & rigid but fragile at interfaces \\
\texttt{cohesion\_integrity\_like} & structural wholeness \\
\texttt{thermal\_transport\_like}  & heat conduction pathway quality \\
\texttt{electrical\_transport\_like} & charge/electron pathway quality \\
\texttt{surface\_reactivity\_like} & surface-dominated chemistry \\
\texttt{fracture\_susceptibility\_like} & failure at weak interfaces \\
\end{tabular}
\end{center}

Do not call them \texttt{youngs\_modulus}, \texttt{hardness}, or
\texttt{conductivity}.  That would be lying in structs.

\subsection{Model Parameters}

\begin{description}
  \item[\texttt{PRECURSOR\_XI\_REF = 20.0}~\AA]
    Reference correlation length for normalisation:
    $\tilde\xi = \clamp(\xi / \xi_\text{ref},\, 0,\, 1)$.
    Represents a characteristic domain size.
  \item[\texttt{PRECURSOR\_MIN\_CONFIDENCE = 0.1}]
    Confidence floor below which a channel is marked invalid.
\end{description}

\subsection{Intermediate Helper Terms}

Three intermediate terms are computed and exposed for traceability before
any precursor channel is evaluated.

\paragraph{Interface penalty.}
\begin{equation}
  P_\text{interface}
    = \alpha\,|\Delta\eta| + \beta\,|\Delta\rho|
    + \gamma\,\Sigma_\text{surface}
  \label{eq:interface_penalty}
\end{equation}
where $\Delta$ are normalised bulk-edge gaps and
$\alpha{=}0.40$, $\beta{=}0.35$, $\gamma{=}0.25$.
Returns~0 when gap fields are NaN (degenerate classifier).
Clamped to $[0,1]$.

\paragraph{Normalised correlation length.}
\begin{equation}
  \tilde\xi = \clamp\!\left(\frac{\xi}{\xi_\text{ref}},\; 0,\; 1\right)
\end{equation}
Returns~0 if $\xi$ is NaN (uniform field or fit failure).

\paragraph{Adaptation capacity.}
When the system is converged (via \fn{compute\_proxy\_delta}),
$A_\text{adapt} = S$ (stabilisation proxy).
Otherwise, $A_\text{adapt} = \bar\eta \cdot (1 - \overline{|\eta-f|})$.
Clamped to $[0,1]$.

\paragraph{Texture lock.}
\begin{equation}
  T_\text{lock} = \bar{P}_{2,\hat{}} \cdot (1 - \min(2\sigma_{P_2},\; 1))
\end{equation}
Orientationally locked structure: high mean alignment with low variance.
Inhibits ductile deformation modes.

\paragraph{Anisotropy index.}
\begin{equation}
  A = \frac{\bar{P}_{2,\hat{}} - 1/3}{1 - 1/3}
      \cdot (0.5 + 0.5\,\tilde\xi)
\end{equation}
Excess texture above the isotropic baseline scaled by correlation
persistence.  Currently isotropic-only; architecture placeholder for
directional $\xi_x, \xi_y, \xi_z$ when available.

\subsection{Precursor Channel Formulas}
\label{sec:precursor_formulas}

All channel values are clamped to $[0, 1]$.
Weights are named constants in \texttt{namespace weights}.

\subsubsection{Rigidity-like}
\begin{equation}
  R^* = w_1\,C + w_2\,U + w_3\,S + w_4\,\tilde\xi
      - w_5\,\Sigma_\text{surface} - w_6\,P_\text{interface}
\end{equation}
where $C$ = cohesion, $U$ = uniformity, $S$ = stabilisation,
$\Sigma$ = surface sensitivity.
Weights: $w_1{=}0.25$, $w_2{=}0.20$, $w_3{=}0.20$, $w_4{=}0.15$,
$w_5{=}0.10$, $w_6{=}0.10$.

\subsubsection{Brittleness-like}
\begin{equation}
  B^* = a_1\,R^* + a_2\,P_\text{interface}
      + a_3\,\Sigma_\text{surface} - a_4\,A_\text{adapt}
\end{equation}
Weights: $a_1{=}0.30$, $a_2{=}0.25$, $a_3{=}0.20$, $a_4{=}0.25$.

\subsubsection{Ductility-like}
\begin{equation}
  D^* = b_1\,U + b_2\,A_\text{adapt} + b_3\,(1 - B^*) - b_4\,T_\text{lock}
\end{equation}
Weights: $b_1{=}0.30$, $b_2{=}0.25$, $b_3{=}0.25$, $b_4{=}0.20$.

\subsubsection{Cohesion integrity-like}
\begin{equation}
  CI^* = f_1\,C + f_2\,S + f_3\,U
\end{equation}
Weights: $f_1{=}0.40$, $f_2{=}0.30$, $f_3{=}0.30$.

\subsubsection{Thermal transport-like}
\begin{equation}
  K^*_\text{th} = c_1\,\tilde\xi + c_2\,T + c_3\,C - c_4\,P_\text{interface}
\end{equation}
Weights: $c_1{=}0.30$, $c_2{=}0.25$, $c_3{=}0.25$, $c_4{=}0.20$.

\subsubsection{Electrical transport-like}
Same form as thermal but heavier weighting on alignment continuity
($c_2{=}0.35$, $c_1{=}0.25$) and equal penalty for interface
discontinuity.

\subsubsection{Surface reactivity-like}
\begin{equation}
  S_r^* = d_1\,\Sigma_\text{surface} + d_2\,P_\text{interface}
        + d_3\,(1 - S) + d_4\,(1 - \tilde\xi)
\end{equation}
Weights: $d_1{=}0.30$, $d_2{=}0.25$, $d_3{=}0.25$, $d_4{=}0.20$.

\subsubsection{Fracture susceptibility-like}
\begin{equation}
  F^* = e_1\,B^* + e_2\,\Sigma_\text{surface}
      + e_3\,P_\text{interface} - e_4\,D^*
\end{equation}
Weights: $e_1{=}0.30$, $e_2{=}0.25$, $e_3{=}0.25$, $e_4{=}0.20$.

\subsection{Confidence Model}

Base confidence starts at~1.0 and is degraded by:
\begin{itemize}
  \item $N < 2 \times \texttt{PROXY\_MIN\_BEADS}$: $-0.3$
  \item Non-finite upstream proxies: $-0.3$
  \item Both $\xi_\eta$ and $\xi_\rho$ NaN: $-0.2$
  \item Actively diverging state ($\Delta\overline{|\eta-f|} > 0.01$): $-0.15$
  \item Degenerate bulk/edge classifier (NaN gaps): $-0.1$
  \item Invalid upstream proxy: confidence $= 0$
\end{itemize}

Individual channels further degrade their confidence:
\begin{itemize}
  \item Transport channels: $\times 0.7$ when no $\xi$ available
  \item Electrical transport: $\times 0.8$ when texture $< 0.01$
  \item Surface reactivity: $\times 0.6$ when no edge beads
  \item Brittleness, ductility, fracture: $\times 0.85$ when no
        delta information (single snapshot)
\end{itemize}

A channel is marked \texttt{valid = false} when its confidence falls
below \texttt{PRECURSOR\_MIN\_CONFIDENCE}.

\subsection{Provenance and Traceability}

Every \type{MacroPrecursorState} records:
\begin{itemize}
  \item \texttt{source\_bead\_count}: $N$ from the upstream proxy
  \item \texttt{source\_valid}: upstream validity flag
  \item \texttt{source\_converged}: upstream convergence flag
  \item \texttt{xi\_norm}, \texttt{interface\_penalty},
        \texttt{adapt\_capacity}, \texttt{texture\_lock}: all
        intermediate terms exposed for inspection
  \item \texttt{anisotropy\_index}: global orientational measure
  \item Per-channel \texttt{confidence} and \texttt{valid}: individual
        trustworthiness
\end{itemize}

Weights are named \texttt{constexpr} constants in
\texttt{namespace weights}, not buried in formulas.

\subsection{Future Calibration Path}

Later versions will fit these channels against:
\begin{enumerate}
  \item Experimental property tables (Young's modulus, hardness,
        conductivity from material databases)
  \item Atomistic simulation benchmarks (MD-computed elastic
        constants, thermal conductivities)
  \item Material family reference sets (metals, polymers, ceramics
        with known property ranges)
\end{enumerate}

The current formulas are explicit heuristic combinations designed for
\textbf{monotone directional correctness}: each perturbation test
verifies the right direction of change, not the absolute magnitude.

\subsection{Directed Precursor Validation (Suite 6B)}

One monotonic perturbation test per precursor channel:

\begin{center}
\begin{tabular}{lll}
\textbf{Channel} & \textbf{Perturbation} & \textbf{Pass criterion} \\
\hline
Rigidity & Dense vs sparse lattice & dense $>$ sparse \\
Fracture susc. & Degrade edge beads & degraded $>$ base \\
Surface react. & Sparse cloud vs dense lattice & surface $>$ bulk \\
Transport & Aligned vs random beads & aligned $>$ random \\
Brittleness & Converged vs fresh & converged $\leq$ fresh $+ 0.05$ \\
Ductility & Lattice vs random cloud & lattice $\geq$ cloud $- 0.1$ \\
Cohesion integ. & Spacing sweep & monotone decrease \\
\end{tabular}
\end{center}

All 85 tests in Suite~\#6 pass (Phases 6A--6E).

% =============================================================================
\section{Property Learning Pipeline}
\label{sec:property_pipeline}
% =============================================================================

This section documents the \textbf{Property Learning Pipeline}
(\file{coarse\_grain/analysis/property\_pipeline.hpp}).  It is the
supervised learning system that maps proxy distributions and precursor
channels to property-scale outputs, with explicit uncertainty
quantification and full provenance.

Architecture position:
\begin{center}
\begin{tabular}{c}
Environment-state evaluation \\
$\downarrow$ \\
Ensemble statistics / spatial field summaries \\
$\downarrow$ \\
Macroscopic response proxies $\to$ precursors \\
$\downarrow$ \\
Property learning pipeline \quad $\leftarrow$ \textit{this layer} \\
$\downarrow$ \\
Calibrated property predictions
\end{tabular}
\end{center}

\subsection{Three-Layer Separation}

The pipeline enforces a strict three-layer architecture.  Formula fusion
is prohibited.

\begin{description}
  \item[Descriptor layer] Computes proxy summaries and precursor channels
    from bead-resolved simulation data.  Contains \textbf{no learned
    parameters}.  Tested in Suites~5--6.
  \item[Calibration layer] Fits model parameters from the exported
    dataset.  Contains \textbf{no simulation logic}.
  \item[Inference layer] Applies the fitted model to new precursor states
    at prediction time.  Contains \textbf{no fitting and no simulation}.
\end{description}

These layers communicate through defined data structures only.  A system
in which physics and learning are fused into a single formula cannot be
audited, cannot be incrementally improved, and cannot be trusted.

\subsection{Model Parameters}

\begin{description}
  \item[\texttt{PIPELINE\_DEFAULT\_REGULARIZATION = 0.01}]
    Default L2 penalty $\lambda$ for Tier~1 ridge regression.
    Prevents overfitting on small datasets.
  \item[\texttt{PIPELINE\_CONFIDENCE\_THRESHOLD = 0.2}]
    Below this, predictions are withheld rather than returned with
    low quality.
  \item[\texttt{PIPELINE\_SINGULAR\_THRESHOLD = 1e-14}]
    Pivot threshold for Gaussian elimination in the linear solver.
  \item[\texttt{PIPELINE\_MAX\_FEATURES = 256}]
    Upper bound on feature vector dimensionality.
  \item[\texttt{ORDINAL\_LOW\_THRESHOLD = 0.33}]
    Boundary for ordinal label assignment: $< 0.33 \to$ Low.
  \item[\texttt{ORDINAL\_HIGH\_THRESHOLD = 0.67}]
    Boundary for ordinal label assignment: $\geq 0.67 \to$ High.
\end{description}

\subsection{Phase 7A: Dataset Construction}

The fundamental unit is \type{PropertyDatasetRow}: a self-contained
record encoding a single simulated material state with five contiguous
blocks:

\begin{center}
\begin{tabular}{cll}
\toprule
\textbf{Block} & \textbf{Content} & \textbf{Struct} \\
\midrule
A & Identity and provenance & \type{DatasetProvenance} \\
B & Full \type{EnsembleProxySummary} & (inline) \\
C & Evaluated \type{MacroPrecursorState} & (inline) \\
D & Property targets (named, typed) & \type{PropertyTarget} \\
E & Training metadata & \type{TrainingMetadata} \\
\bottomrule
\end{tabular}
\end{center}

\paragraph{Target type system.}
Each target carries a \type{PropertySourceType} classification:
\texttt{Experimental}, \texttt{Atomistic}, \texttt{Synthetic}, or
\texttt{Ordinal}.  Missing targets are represented with
\texttt{available = false} and are \textbf{never} silently substituted
with zero.

\paragraph{Row validity.}
A row is valid when: (1)~\texttt{system\_id} is non-empty,
(2)~at least one target is available, and
(3)~\texttt{source\_bead\_count} $> 0$.

\subsection{Phase 7B: Feature Vectorization}

\type{ModelFeatureConfig} specifies which feature blocks to include.
\fn{vectorize()} converts a dataset row to a fixed-length
\type{FeatureVector} with deterministic column ordering.

\begin{center}
\begin{tabular}{clc}
\toprule
\textbf{Block} & \textbf{Content} & \textbf{Count} \\
\midrule
1 & Proxy scalars + reference-normalised forms & 10 \\
2 & Structural modifiers (penalty, gaps, $\tilde\xi$) & 7 \\
3 & Precursor channel values & 8 \\
4 & Confidence and support terms & 6 \\
5 & Categorical context (deferred) & --- \\
\midrule
  & Total (default config) & 31 \\
\bottomrule
\end{tabular}
\end{center}

\paragraph{Missing value policy.}
Three policies: \texttt{Zero} (replace with 0, counted as imputed),
\texttt{NaN} (pass through, not imputed), \texttt{Mean} (fallback to 0
when no mean vector available).

\paragraph{Ablation.}
Each block can be independently disabled via \type{ModelFeatureConfig}
boolean flags, enabling systematic ablation studies.

\subsection{Phase 7C: Model Architecture}

\subsubsection{Tier 1: Transparent Linear Baseline}

Regularized ridge regression solves:
\begin{equation}
  \mathbf{w} = (\mathbf{X}^\top \mathbf{X} + \lambda \mathbf{I})^{-1}
               \mathbf{X}^\top \mathbf{y}
  \label{eq:ridge}
\end{equation}

The linear solver uses Gaussian elimination with partial pivoting.
Coefficients are stored in \type{LinearModelCoefficients} with named
feature lookup and explicit residual standard deviation.

\paragraph{Prediction output.}
Every prediction is returned as \type{PropertyPrediction}:
\begin{itemize}
  \item \texttt{value}: predicted scalar
  \item \texttt{uncertainty}: residual standard deviation from training
  \item \texttt{confidence}: derived from input validity, imputation count,
        and upstream convergence quality
  \item \texttt{top\_features}: ranked contributing features via
        $|w_i \cdot x_i|$
  \item \texttt{withheld}: \texttt{true} when confidence $<$ threshold
\end{itemize}

\paragraph{Abstention.}
The system does not produce confident predictions from poorly constrained
inputs.  When confidence falls below
\texttt{PIPELINE\_CONFIDENCE\_THRESHOLD}, the prediction is withheld.

\subsubsection{Tier 2 and Tier 3}

Tier~2 (gradient boosted trees, shallow MLP) and Tier~3
(family-conditioned model heads) are specified in the architecture but
deferred to later implementation.  They share the same feature interface
and prediction output structure.

\subsection{Phase 7D: Evaluation and Split Discipline}

\subsubsection{Metrics}

\begin{center}
\begin{tabular}{ll}
\toprule
\textbf{Task Type} & \textbf{Primary Metrics} \\
\midrule
Regression & RMSE, MAE, Spearman $\rho$ \\
Ranking & Pairwise accuracy, Kendall $\tau$ \\
\bottomrule
\end{tabular}
\end{center}

Spearman rank correlation uses rank-transformed Pearson correlation.
Kendall $\tau$ counts concordant vs.\ discordant pairs.

\subsubsection{Grouped splits}

Rows are partitioned by \texttt{split\_group}, ensuring all perturbed
variants of a single base structure remain in the same partition.
\fn{verify\_split\_discipline()} confirms no group crosses partition
boundaries.  This prevents the model from learning family fingerprints
rather than generalizable descriptor-to-property mappings.

\subsection{Phase 7E: Synthetic Supervision}

\subsubsection{Training stages}

\begin{enumerate}
  \item \textbf{Stage~1 (synthetic)}: Controlled perturbation sweeps
    whose ground-truth orderings are known by construction.
    \fn{build\_synthetic\_sweep()} generates dataset rows by
    monotonically varying one proxy field while holding others fixed.
    \fn{verify\_monotone()} validates the predicted ordering.
  \item \textbf{Stage~2 (ordinal)}: Weak labels from domain rules.
    \fn{assign\_ordinal\_label()} maps channel values to
    \{Low, Medium, High\} at thresholds $0.33$ and $0.67$.
  \item \textbf{Stage~3 (external)}: Reference values from measured
    property databases or atomistic simulation outputs.  Deferred
    pending availability of reference datasets.
\end{enumerate}

\subsubsection{Deployment order}

Ranking $\to$ ordinal classification $\to$ regression.  This ordering
reflects the practical difficulty of obtaining reliable numeric property
labels.

\subsection{Directed Pipeline Validation (Suite 7)}

\begin{center}
\begin{tabular}{lll}
\textbf{Phase} & \textbf{Focus} & \textbf{Tests} \\
\hline
7A & Dataset row validity, targets, provenance & 15 \\
7B & Vectorization determinism, ablation, missing policy & 15 \\
7C & Linear fit, coefficients, prediction, abstention, attribution & 15 \\
7D & RMSE/MAE, Spearman, Kendall, grouped splits & 15 \\
7E & Ordinal labels, monotone sweeps, end-to-end pipeline & 49 \\
\hline
   & Total & 109 \\
\end{tabular}
\end{center}

All 109 tests in Suite~\#7 pass (Phases 7A--7E).

% =============================================================================
\section{Property Calibration and Target Legitimacy}
\label{sec:property-calibration}
% =============================================================================

\subsection{Overview}

The \textbf{Property Calibration and Target Legitimacy} layer
(\file{property\_calibration.hpp}) provides a disciplined framework for
mapping precursor channels to a compact, defensible set of property-scale
outputs with explicit calibration logic, confidence decomposition,
rejection behaviour, and a formal promotion policy governing when a proxy
quantity may be named as a material property.

This module operates in namespace \ns{coarse\_grain::calibration} and
depends on the property pipeline (\ns{coarse\_grain::pipeline}), the macro
precursor layer, and the ensemble proxy layer.  It introduces no simulation
logic and contains no learned parameters that bypass the three-layer
separation.

\subsection{Architecture Position}

\begin{verbatim}
Environment-state evaluation
    |
Ensemble statistics / spatial field summaries  (ensemble_proxy.hpp)
    |
Macroscopic response proxies -> precursors      (macro_precursor.hpp)
    |
Property learning pipeline                      (property_pipeline.hpp)
    |
Property calibration and target legitimacy   <- this module
\end{verbatim}

\subsection{Phase 8A: Property Family Selection}

Five first-wave property targets are selected:
\begin{enumerate}
  \item \textbf{rigidity\_like} --- structural resistance to deformation
  \item \textbf{ductility\_like} --- capacity for adaptive, non-fracturing deformation
  \item \textbf{brittleness\_like} --- susceptibility to abrupt failure
  \item \textbf{thermal\_transport\_like} --- heat conduction capacity
  \item \textbf{electrical\_transport\_like} --- electrical conduction capacity
\end{enumerate}

The \texttt{-like} suffix is retained throughout proxy-only and
calibrated-relative states.  It may only be dropped when a target reaches
externally-anchored legitimacy under Phase~8H.

\subsection{Phase 8B: Target-Definition Contracts}

Each property target has a formal \type{TargetContract} specifying:
\begin{itemize}
  \item \textbf{Sign expectations}: which precursor channels increase
        (positive) or decrease (negative) the target
  \item \textbf{Domain constraints}: admissible material classes,
        excluded regimes, scale assumptions
  \item \textbf{Output range}: valid $[\min, \max]$ for predictions
  \item \textbf{Confidence requirements}: minimum precursor confidence,
        minimum bead count, maximum imputed features
  \item \textbf{Dominant features}: for OOD detection and attribution checks
\end{itemize}

The contract is the standard against which both analytic formulae and
trained models are evaluated.  A model that violates the sign expectations
of the contract has failed a scientific test.

\subsection{Phase 8C: Calibration Profiles}

Each property target carries a \type{CalibrationProfile} providing a
calibrated interpretation layer without modifying the descriptor or
precursor computation code:

\begin{lstlisting}
struct CalibrationProfile {
    std::string property_name;
    double scale, offset;
    double min_confidence_threshold;
    double abstention_threshold;
    bool   monotone_expected;
    std::vector<std::string> dominant_features;
    double apply(double raw) const;
};
\end{lstlisting}

Ductility is marked as non-monotone (has an optimum); all other first-wave
targets are marked monotone.

\subsection{Phase 8D: Supervision Regime Taxonomy}

Three supervision regimes are defined:

\begin{enumerate}
  \item \textbf{Contract-supervised}: no external data; sign and monotonicity
        enforced by contracts.  Outputs are physics-informed ranks.
  \item \textbf{Synthetic-calibrated}: controlled perturbation sweeps provide
        relative target labels.  Outputs are synthetic training artifacts.
  \item \textbf{Externally calibrated}: measured or atomistically validated
        reference values anchor outputs.  Only this regime supports promotion
        to named properties.
\end{enumerate}

The regime under which any prediction was produced is recorded in
\type{CalibratedPrediction::regime} and reported alongside the output.

\subsection{Phase 8E: Confidence-Gated Prediction}

Confidence is decomposed into five named components:

\begin{center}
\begin{tabular}{ll}
\textbf{Component} & \textbf{Source} \\
\hline
\field{input\_confidence}       & Feature validity flags and imputation count \\
\field{precursor\_confidence}   & Convergence quality and precursor validity \\
\field{coverage\_confidence}    & OOD envelope check (Phase 8F) \\
\field{attribution\_stability}  & Top contributing features stable \\
\field{calibration\_confidence} & Prediction within valid calibration range \\
\end{tabular}
\end{center}

Total confidence is the product of all five components.  Individual
component floors are enforced independently: a prediction may pass the
total threshold while failing on a single dimension.

Three prediction states are defined:
\begin{itemize}
  \item \textbf{Accepted}: total and all components above thresholds
  \item \textbf{LowConfidence}: total above abstention but at least one
        component marginal
  \item \textbf{Withheld}: total below abstention or hard constraint violated
\end{itemize}

\subsection{Phase 8F: Out-of-Distribution Detection}

A \type{FeatureEnvelope} stores per-feature mean and standard deviation
from the training partition.  For each dominant feature, a $z$-score is
computed against the envelope; features exceeding the configured threshold
(default $z > 3$) are flagged as extrapolating.

Extrapolation sets \field{coverage\_confidence} to zero, which forces the
total confidence below any abstention threshold, producing a Withheld
status.  The extrapolation tag is recorded in prediction provenance.

Only dominant features (as specified by the target contract) trigger
extrapolation tagging.

\subsection{Phase 8G: Property-Specific Synthetic Curricula}

Per-property sweeps are constructed in which relevant precursor channels
are varied monotonically while others are held fixed:

\begin{itemize}
  \item \textbf{Rigidity}: sweep cohesion $0 \to 1$; expect monotone increase
  \item \textbf{Transport}: sweep texture $0 \to 1$; expect monotone increase
  \item \textbf{Brittleness}: sweep surface sensitivity $0 \to 1$ with
        decreasing stabilisation; expect monotone increase
\end{itemize}

All curriculum rows are tagged with \field{SyntheticCalibrated} regime.
A model that fails to recover expected responses on controlled sweeps has
not learned the physics it is intended to represent.

\subsection{Phase 8H: Target Promotion and Legitimacy}

Three legitimacy states are defined:

\begin{center}
\begin{tabular}{lll}
\textbf{State} & \textbf{Suffix} & \textbf{Requirements} \\
\hline
ProxyOnly            & \texttt{-like} retained & No additional evidence \\
CalibratedRelative   & \texttt{-like} retained & Monotonicity, signs, domain documented \\
ExternallyAnchored   & suffix dropped           & External data, bounded error, validated rejection \\
\end{tabular}
\end{center}

Promotion from ProxyOnly to CalibratedRelative requires preserved
monotonicity, consistent sign behaviour, and a documented domain statement.
Promotion from CalibratedRelative to ExternallyAnchored additionally
requires external data coverage, bounded residual error, and validated
rejection behaviour on out-of-domain inputs.

Neither promotion is permanent.  If calibration drift is detected, if
contract violations appear, if coverage confidence degrades, or if
attribution patterns become inconsistent, the target is demoted and the
demotion is recorded in model provenance.

Multi-dimensional grouped splits are provided via
\op{multi\_grouped\_split()}, which groups rows jointly on material family,
generator name, split group, and fragment family to prevent leakage.

\subsection{Testing Summary}

Suite~\#8 validates all eight phases:

\begin{center}
\begin{tabular}{lll}
\textbf{Phase} & \textbf{Focus} & \textbf{Tests} \\
\hline
8A & Property family selection, naming, indexing & 16 \\
8B & Target contracts, sign expectations, domain & 17 \\
8C & Calibration profiles, transform, clamping & 12 \\
8D & Supervision regime taxonomy & 4 \\
8E & Confidence decomposition, status, attribution & 23 \\
8F & OOD detection, envelope, extrapolation & 14 \\
8G & Synthetic curricula, monotonicity, regime tagging & 14 \\
8H & Legitimacy states, promotion, demotion, splits & 93 \\
\hline
   & Total & 193 \\
\end{tabular}
\end{center}

All 193 tests in Suite~\#8 pass (Phases 8A--8H).

\end{document}
